\documentclass{jsarticle}
\usepackage[pdftex]{graphicx}
\usepackage{mathtools,amsmath,amssymb}
\usepackage{multirow}
\usepackage{xcolor}
\usepackage{bm}
\usepackage[dvipdfmx]{hyperref}

\def\be#1\ee{\begin{align}#1\end{align}}


\begin{document}

-----------------------
畑中さん
-----------------------\\

\textcolor{blue}{論文で具体的に扱っているSAWの周波数は4$\sim$5GHzなので、その場合の波長は$\lambda_a=$1umかまたは数100nm程度かと思われます。一方で、論文中の計算で改定しているYIG板の膜厚はL$=$500umとのことですので、膜厚はSAW波長より十分に大きい構造となっています($L\gg\lambda_a$)。この場合では、YIGの表面と裏面でのモードが結合して対称・反対称Lambモードを形成するほど強い結合は得られず、むしろ、ほぼ両面で独立したSAWレイリー波に近い波が生じていると考えた方が自然なのではと感じました。もちろん、$L\gg\lambda_a$であっても、SAWのような半無限体を仮定した固有モードでなく、有限膜厚における固有モードなので、Lamb波の一種だと思いますが、その特性はSAWにほぼ漸近していると思われます。一般的にLamb波を考える上では$L\sim\lambda_a$が妥当な構造ではと感じました。L=数umで計算した場合では、Fig2と似たような結果が得られないでしょうか?もしそうならその結果を見せた方が一般的なLamb波を扱う構造に近いですし、おそらく、モードオーバーラップも大きくなるので、より鮮明なピークの反交差が現れるのではと推察します。}\\

ご指摘ありがとうございます。そして大変申し訳ないのですが、本文中の単位を間違っておりました。計算は$L=500$ nmの系で実施しています。前回の議論で実験的には$200$ nmくらい以下では脆くて作りづらく、$1\ \mu$mくらいまでしか薄膜っぽくない、と伺っていたのでその中間くらいの厚さにしたのでした。なのですでにおおよそ$L\sim\lambda_a$くらいの設定になっているかと思います。こんなつまらないミスのせいでお時間を浪費してしまい本当に申し訳ないです...\\

\textcolor{blue}{P. 4の左側の第4パラグラフで、”… BVWs appear, plotted as blue and red curves …”と記載されていますが、Fig. 2(a)のBVWの分散曲線の青と赤が見えにくいです。黒点をもう少し小さくして、赤・青線の視認性を上げることは可能でしょうか。}\\

ご指摘ありがとうございます。ズームイン以外はマーカーを中抜きにして線幅を太くしてみました。これで見やすくなりましたでしょうか。\\

\textcolor{blue}{Fig. 2aで、BVWの基底モードである最低周波数は、膜厚方向に一様なモード分布形状で、symmetricな対称モードになるかおもいました(勘違いでしたら申し訳ございません)。ただ、計算では赤線のasymmetricな反対称モードになっているようです。こちらは膜厚方向でどのようなモード形状を持っていますか?}\\

ご指摘ありがとうございます。結論だけ先に申し上げますと、BVW, Lamb waveともに「反対称モード」が基底モードで間違いないようです。ただ、この呼称はBVWについては適切ではなさそうなので、本文では修正しました。お恥ずかしくも私自身の理解がだいぶあやふやでしたので、以下に内容をまとめさせていただきます。\\

そもそも、symmetric mode, anti-symmetric modeという呼称はLamb waveの方から来ているので、そちらからまとめます。Lamb (1917)~\cite{lambWavesElasticPlate1917}の記法を我々のものと合わせると、
\be
    \mathrm{Symmetrical:}&\ u_x(-z)=u_x(z),\quad u_z(-z)=-u_z(z),\\
    \mathrm{Asymmetrical:}&\ u_x(-z)=-u_x(z),\,u_z(-z)=u_z(z).
\ee
として定義されています。これは、$z=0$について、膜厚方向の変位$u_z$が対称(odd func.)か、反対称(even func.)かに対応しています。我々の議論している$C_{2x}$-rotational symmetry については、$\bm{u}$が固有状態になっている必要があるので、$C_{2x}\bm{u}(z)=\lambda\bm{u}(z)$とすると、$C_{2x}\bm{u}(z)=(u_x(-z),-u_y(-z),-u_z(-z))$より、
\be
    \mathrm{Symmetrical:}&\ C_{2x}\bm{u}(z)=+\bm{u}(z),\\
    \mathrm{Asymmetrical:}&\ C_{2x}\bm{u}(z)=-\bm{u}(z),
    \ee
と(a)symmetrical modeと固有値が対応します。\\
    
この対応をまねて、BVWについてもsymmetric, antisymmetric modeを
\be
    \mathrm{Symmetric:}&\ C_{2x}\bm{m}(z)=+\bm{m}(z),\\
    \mathrm{Antisymmetric:}&\ C_{2x}\bm{m}(z)=-\bm{m}(z), 
\ee
として定義していました。今度は振動は$yz$面内で起こり、$C_{2x}\bm{m}(z)=(m_x(-z),-m_y(-z),-m_z(-z))$であるので、
\be
    \mathrm{Symmetric:}&\ m_{y,z}(-z)=-m_{y,z}(z),\\
    \mathrm{Antisymmetric:}&\ m_{y,z}(-z)=+m_{y,z}(z),
\ee
となり、成分で見ると奇関数がsymmetric, 偶関数がanti-symmetricになっています。つまり、基底モードはanti-symmetric modeであり、磁化の振動成分は$z$について偶関数的なモード形状になっています。\\

このように、$C_{2x}$固有値によって名前を付けてしまうと、各成分が$z$について偶関数か奇関数かが直感的に(?)一致しなくてややこしいので、``We define the angle...''の段落を書き換え、$C_{2x}$固有値によってグラフの色を付けている。ということにしました。\\

-----------------------
山本さん（一部）
-----------------------\\

---YIGのパラメータ導入少し上、generalized eigenvalue problemについて---\\
\textcolor{blue}{固有値は小さい方から順番に有限個計算したのでしょうか?exchangeを入れていないので有限の周波数区間に無限個のスピン波モードがある状態になっており、思ったモードに収束していない可能性があります。}\\

申し訳ありません。危惧されている点を理解しきれなかったので数値計算で行っていることを報告します。\\
やっていることとしては$z$方向の離散化です。$N$個の小区間に分割し、有限要素法を用いて形状関数を基底に固有関数を数値近似しています。この際、Navier--Cauchyは時間について二階の微分方程式なので$(u_x,u_z)$に加え、$(\dot{u}_x,\dot{u}_z)$も基底に加えています。これらと磁化の$(m_2,m_3)$の二成分(2,3は面内磁場に垂直で互いに垂直な方向)を加えて、$6N\times 6N$の行列の対角化を各$k$に対して行っています。\\

Lamb modeは$4N$本出てきますが、一番下の二つ以外は画面外の上にいます。BVWは$2N$本出てきています。$N=200$で計算し、BVWの方は全てのモードをプロットしています。上界となる周波数は解析的な値$(\omega=\gamma\mu_0\sqrt{H_0(H_0+M_0)}\approx 4.96 \mathrm{\ GHz})$に一致しています。\\

---$\phi=\pi/2$の場合の対称性の議論について---\\
\textcolor{blue}{この主張は真でしょうか?ということはLamb波ではなくてshear horizontalモードはスピン波と結合するということですか?}\\

たぶん真で、おっしゃる通りhear horizonta modeは結合するのではないかなと思います。まず、段落の議論の詳細を述べます。この場合、$\bm{M}_0=(0,M_{0,y},0)$であり、鏡映操作$M_y$によって$\bm{M}_0\to(-M_{0,x},M_{0,y},-M_{0,z})$となることから、系は$y$軸に垂直な鏡映面($xz$平面)についての鏡映対称性を持ちます。Lamb波は$(u_x,u_z)$の振動なので、$M_y\bm{u}=+\bm{u}$を満たし、鏡映固有値は$+1$です(symmetric)。磁化についても、$y$軸周りに歳差するので振動は$(m_x,m_z)$で起こります。したがって、$M_y$によって反転し、$M_y\bm{m}=-\bm{m}$となります。これら$\bm{u}$と$\bm{m}$は規約表現が異なるので結合は起こらないはずです。\\

次に、ご指摘のShear horizontal modeについて、こちらは$u_y$の振動となり、symmetricもanti-symmetricも生じうるので、anti-symmetricなmagnetostatic surface waveとcoupleするモードも生じると考えられます。\\


-----------------------
浅野さん
-----------------------\\

\textcolor{blue}{この理論構築の新規部分をもう少し主張した方が良いように感じます。Kalnikos-Slavinが与えたダイポール相互作用に関するグリーン関数を変位場で摂動展開し、変位場とマグノンに対する結合方程式を厚み方向に対して等価に扱った点が新規という理解でしょうか?あるいはそれを数値的に解いた部分でしょうか?既存理論枠組みで議論されてきた論文を引用しながら、今回の理論構築で新たに予言されること、新たな適用範囲などをイントロ部などで整理したほうが良いかと思いました。}\\

ありがとうございます。本論文の新規性としては、静磁波、Lamb波それぞれは詳細に議論されてきたが、双極子相互作用を中心的に扱ってそれらの結合を等価に扱った理論は存在しなかったのでそれを作った。というところかと思います。結果として双極子相互作用によるmode mixingの大きさを定量的に見積もれるようになったという点が新たに予言できるようになったことだと思います。イントロ部分を書き直しました。\\

\textcolor{blue}{数値結果で対称性から結合の様子を議論するのはよくわかるのですが、分裂幅several MHzの正当性を議論する必要は無いでしょうか?マグノメカニクスの潮流的にこのあたり実験屋さんが一番気になるところかと思います。数GHzの共振器でQ値1000くらいを仮定するとラフには強結合領域ですかね。}\\

ありがとうございます。結合定数周辺の話や実験的に見れるかどうかなどの具体的なところは完全に勉強不足だったのですが、conclusionの前に一段落入れてgapサイズについて議論してみました。記述がおかしくないか確認していただけますでしょうか。\\


\textcolor{blue}{45度の結果で中央下部のインセットですが、マグノン側は2本の線がほとんど重なって交差するように見受けられますが、これはそのうちの1本との交差を示しているということでしょうか。あとFig.2の横軸のkLのラベルが全体的に右側にずれているように見受けられますが中央に表示した方が良いかと思います(kL$=$1の部分に書いている?)。}\\

はい。まず、マグノン側は基底モード1本との交点をズームしてみています。このすぐ上に数十本のモードがいる感じです。Fig.~2のcaptionの最後に説明を追加しました。\\

ラベルについてご指摘ありがとうございます。中央にずらすように変更しました。今まであまり意識せず、目盛りの無い空いている空間に適当にラベルを入れ込んでいました。\\

\textcolor{blue}{ひょっとすると今回の論文から敢えて外されているのかもしれませんが、以前の打ち合わせでコメントさせていただいた通り、既存理論の磁気弾性定数と今回の理論で算出されるものの対応関係は調べたほうが良いかと思いますがいかがでしょうか。}\\

すみません。今回の論文からは意図的に外させていただきました。理由としては、前回議論させていただいた際には磁気弾性定数と比較できる。と申し上げたのですが、今の計算の次数では出せない気がしてきたためです。というのも、通常の磁気弾性結合のマクロな理論では$\epsilon_{ij}m_im_j$までで自由エネルギーを求め、そこから汎関数微分で磁場を出すという形かと思います。これにより$m$と$\epsilon (\sim u)$の交差項の係数が分かるという流れです。ただ、今回は有効磁場を$m,u$の一次までしか計算していないので磁気弾性定数の定義式より次数が低くなっていて、更に高次の展開をしないと既存理論との比較はできないのではないかと考えました。\\

高次展開は原理的には可能なので今後の課題というか、次にまとめて出す論文の方に譲り、今回の論文では$m,u$の最低次での基本方程式の導出とバンドの再現に主眼を置くという形にしたいのですがいかがでしょうか。\\

\textcolor{blue}{あまり本質的ではありませんが、L$=$500umでThin filmというのは少し違和感がありました。}\\

すみません。こちら私のミスで、計算は$500$nm厚の系で行っていました。このようなことでお時間を取ってしまい申し訳ないです...\\

\textcolor{blue}{完全に数式を追い切れておりませんで恐縮ですが、磁気弾性結合によるフォースfに近似なしでikが出てくるのはLamb波だからという理解で正しいでしょうか。通常のGivenな磁気弾性結合$\epsilon_{ij} x_i x_j$から出すときはフォノンの分布に応じてGeneralな微分項(例えば$\partial u_x/\partial z$など)が出てきて、近似的にikに比例する項だけ残す営みかと思いましたので質問させていただきました。}\\

これはLamb波の特性からというよりは、$x$方向に平面波であることを最初に仮定しているから出てくるのだと思います。独立な波の形として$e^{ikx}$に比例することを仮定しているので、すべての計算過程で$\partial_x\to ik$と置き換えています。一方で、定式化でも数値計算でも厚み方向の微分 $\partial_z$ は残したまま計算しています。


\begin{thebibliography}{99}
\bibitem{lambWavesElasticPlate1917}
H.~Lamb, On Waves in an Elastic Plate, \href{https://doi.org/10.1098/rspa.1917.0008}{Proc. A \textbf{93}, 114 (1917)}.

\end{thebibliography}

\end{document}